\uuid{8lDn}
\exo7id{3420}
\titre{exo7 3420}
\auteur{quercia}
\organisation{exo7}
\datecreate{2010-03-10}
\isIndication{false}
\isCorrection{false}
\chapitre{Déterminant, système linéaire}
\sousChapitre{Système linéaire, rang}
\cpge{mpsi-mp2i.applications-lineaires.e, pcsi.applications-lineaires.f, mpsi-mp2i.matrices-applications-lineaires.c, pcsi.matrices-applications-lineaires.d, ptsi.matrices-applications-lineaires.d, bcpst1.systemes-lineaires.a, bcpst2.applications-lineaires-matrices.a, bcpst2.applications-lineaires-matrices.b}
\module{Algèbre}
\niveau{L2}
\difficulte{}

\contenu{
\question{
Soient $f,f_1,\dots,f_p$ des formes linéaires sur $ K^n$ linéairement
indépendantes.
Montrer que $f$ est combinaison linéaire de $f_1, \dots, f_p$ si et seulement si
$\mathrm{Ker} f \supset \mathrm{Ker} f_1 \cap \dots \cap \mathrm{Ker} f_p$.

Indication :
\'Etudier le système $$ \left\{\begin{array}{lll} f_1(x) &=& 0 \cr
                               \quad\vdots  \cr
                               f_p(x) &=& 0 \cr
                               f(x)   &=& 1.\cr \end{array}\right.$$


Le résultat est-il encore vrai si on ne suppose pas $(f_1, \dots, f_p)$
libre ?
}
}
